% !TEX TS-program = xelatex
% !TEX encoding = UTF-8 Unicode
% !Mode:: "TeX:UTF-8"
% Tailored for: Biopharma R&D / QC / Analytical roles

\documentclass{my-resume}
\usepackage{my-resume-zh}
\usepackage{linespacing_fix}
\usepackage{cite}

\begin{document}
\pagenumbering{gobble}

\photoheader{%
  \name{吴筱琦}
  \contactInfo{+86 13776020628}{yukiwxq@gmail.com}{现居城市：伦敦/苏州}{意向城市：上海/苏州/杭州 \textbar\ 生物医药研发/QC/分析岗}
}{../template/photo.jpg}{2.4cm}

\section{教育背景}
\entry{帝国理工学院 (QS 2)}{2025.09 - 2026.09}{计算生命科学研究型硕士}{伦敦}
\entry{帝国理工学院 (QS 2)}{2022.09 - 2025.07}{生物科学 本科}{伦敦}
\begin{itemize}[parsep=0.2ex]
  \item \textbf{相关课程}：应用分子生物学，遗传学，细胞生物学与遗传学，生物化学与微生物学，生物信息学，编程与统计
\end{itemize}

\section{实习经历}
\entry{上海宣泰医药科技有限公司}{2024.07 - 2024.09}{分析工程师 \textperiodcentered\ 分析部}{上海}
\begin{itemize}[parsep=0.2ex]
  \item \textbf{稳定性检测}：使用高效液相色谱-紫外检测（HPLC/UV）技术对客户提供的药品样本进行稳定性检测，专注于样品中杂质的检测与分析
  \item \textbf{数据分析}：独立完成样本前处理、HPLC系统操作及数据分析，确保检测结果的准确性和可靠性
  \item \textbf{质量控制}：通过严格的质控流程鉴别潜在杂质，为药品的稳定性评估提供关键数据支持
\end{itemize}

\section{科研与项目经历}

\entry{多模态大语言模型提取媒介性状数据评估：高精确率与结构依赖型记录召回}{2025.12 - 2026.08}{}{伦敦}
\role{Python、Gemma 4 31B IT（多模态LLM）、PyMuPDF、Docling、提示工程、JSON Schema 校验、pytest}{}
\begin{itemize}[parsep=0.2ex]
  \item \textbf{基准构建与提取流程}：设计加权多样性贪心算法从 VecTraits 超4万行、334篇论文中选取20篇，构建2,335条开发基准与431条留出测试集；结合 PyMuPDF/Docling 证据提取，驱动 Gemma 4 31B IT 完成提示工程约束下的可溯源性状记录提取
  \item \textbf{严谨评估与质量验证}：设计二分图记录匹配评估框架，在留出集上验证锁定流程，取得记录精确率 1.000、条件宏观字段 F1 0.729，体现与药品质量控制一致的系统性数据验证思维
\end{itemize}

\entry{病毒指纹图谱：利用 HERV-K solo-LTR 多态性进行人群遗传学分析}{2023.11 - 2023.12}{}{伦敦}
\role{PCR、凝胶电泳、群体遗传学分析}{}
\begin{itemize}[parsep=0.2ex]
  \item \textbf{实验操作}：运用 PCR 和凝胶电泳检测 5 组引物对应的基因组多态性，结合 F$_{ST}$（范围 0.0010-0.3659）和 HWE（$\chi^2$ 检验，临界值 3.841）进行群体遗传学分析
  \item \textbf{分析评估}：评估 HERV-K solo-LTR 作为人群遗传与法医学标记的潜力与挑战，识别显著偏离 HWE 的引物位点
  \item \textbf{研究成果}：完成实验设计、操作与数据分析全流程，成果获优秀评价
\end{itemize}

\entry{OCA1A 型白化病的生物信息学分析}{2023.10 - 2023.11}{}{伦敦}
\role{NCBI、OMIM、Ensembl、UniProt、AlphaFold、ClinVar}{}
\begin{itemize}[parsep=0.2ex]
  \item \textbf{结构分析}：分析 TYR 基因突变对酪氨酸酶结构与功能的影响，聚焦铜结合域突变，并使用 AlphaFold 结构预测辅助分析
  \item \textbf{数据库整合}：跨 NCBI、OMIM、Ensembl、UniProt、ClinVar 数据库进行变异注释与人群频率分析
  \item \textbf{研究成果}：完成生物信息学报告，掌握多数据库整合、蛋白结构预测与进化保守性分析技能
\end{itemize}

\entry{近交系小鼠在医学研究中的应用与局限性综述}{2023.11 - 2023.12}{}{伦敦}
\role{文献综述、遗传学分析}{}
\begin{itemize}[parsep=0.2ex]
  \item \textbf{证据综合与批判分析}：阐述近交系小鼠遗传学原理及其临床前药物应用，并批判分析其遗传/生理差异等局限性，成果获优秀评价
\end{itemize}

\section{技能/其他}
\entrySkillsSep
\begin{itemize}[parsep=0.2ex]
  \item \textbf{分子与生化实验}：PCR 与凝胶电泳、DNA/RNA 提取（Qiagen 试剂盒）、蛋白质电泳、西方印迹；群体遗传学分析（F$_{ST}$、Hardy-Weinberg 平衡/$\chi^2$ 检验）；跨 NCBI、OMIM、Ensembl、UniProt、ClinVar、AlphaFold 的变异注释与蛋白结构分析
  \item \textbf{数据分析}：R（ggplot2、dplyr）、Python、统计推断（GLM、Wilcoxon 检验、t 检验、Pearson 相关）、质量控制文档规范
  \item \textbf{语言}：英语（雅思 7.5），普通话（母语），法语（A2）；\textbf{兴趣爱好}：钢琴，滑雪
\end{itemize}

\end{document}
